\begin{table}
    \footnotesize
    \centering
    \begin{threeparttable}
        \caption{Modal Verb Strength of authors' \(t\)th paper (draft and final)}
        \label{table9llmmodalverb}
        \begin{tabular}{p{3cm}S@{}S@{}S@{}S@{}S@{}S@{}}
            \toprule
            &{\(t=1\)}&{\(t=2\)}&{\(t=3\)}&{\(t=4\text{--}5\)}&{\(t\ge6\)}\\
            \midrule
            \multicolumn{6}{l}{\textbf{Predicted \(R_{jP}-R_{jW}\)}}\\
            \quad Women                   &       -0.27   &       -0.11   &        0.16   &        0.25   &        0.44   \\
                                          &      (0.34)   &      (0.24)   &      (0.28)   &      (0.41)   &      (0.56)   \\
            \quad Men                     &        0.31***&        0.30***&        0.38***&        0.30***&        0.32***\\
                                          &      (0.08)   &      (0.05)   &      (0.03)   &      (0.05)   &      (0.08)   \\
            \midrule\multicolumn{6}{l}{\textbf{Marginal effect of female ratio}}\\
            \quad Published article       &       -0.56   &       -0.43   &       -0.30   &       -0.16   &       -0.03   \\
                                          &      (0.66)   &      (0.48)   &      (0.36)   &      (0.37)   &      (0.50)   \\
            \quad Draft paper             &        0.02   &       -0.03   &       -0.07   &       -0.11   &       -0.15   \\
                                          &      (0.52)   &      (0.39)   &      (0.31)   &      (0.29)   &      (0.35)   \\
            \midrule
            \textbf{Diff.-in-diff.}&       -0.58   &       -0.40   &       -0.23   &       -0.05   &        0.12   \\
                                          &      (0.39)   &      (0.28)   &      (0.30)   &      (0.43)   &      (0.61)   \\
            \bottomrule
        \end{tabular}
        \begin{tablenotes}
            \tiny
            \item \textit{Notes}. Sample 4,289 observations. Panel one displays magnitude of predicted \(R_{jP}-R_{jW}\) (the direct effect of peer review) for women and men over increasing \(t\). Panel two estimates the marginal effect of an article's female ratio (\(\beta_1+\beta_2\times t\)), separately for draft papers and published articles. Figures from FGLS estimation of~equation~(15), weighted by \(N_{it}\) (see~the data appendix). Control variables include citation count (asinh), \(\text{max. }T\) (author prominence) and \(\text{max. }t\) (author seniority), native speaker and editor and journal-year fixed effects. Standard errors clustered by editor and robust to cross-model correlation in parentheses. ***, ** and * statistically significant at 1\%, 5\% and 10\%, respectively.
        \end{tablenotes}
    \end{threeparttable}
\end{table}